\section{模型假设}

\begin{enumerate}
  \item 将药材视为连续介质，以有效热物性和扩散系数描述其内部传递过程，不单独分辨细胞、孔隙等微观结构。
  \item 以附件给出的烘房温度与水分浓度代表药材侧壁周围的环境状态，忽略侧壁不同位置的环境差异，表面交换由题目给定的对流参数描述。
  \item 不考虑药材内部化学反应等过程产生的体热源。
  \item 烘干阶段切换时，药材内部温度与含水率连续变化，不因环境边界改变而人为重置。
\end{enumerate}

固定几何、常物性及潜热处理等局部选择在对应小问中说明。题目给定的尺寸、初值和物性参数作为已知条件使用。
